\cotefile{71}
\scanpage{36}\mspage{$\underline{\overline{\mathrm{IX}}}$}
Comme $\Gamma$ est un $J$-polygone et $\partial E(\Gamma)$ est un polygone convexe il est possible de trouver un disque $D_s$ qui vérifie b) de la def.~1 par rapport à $\Gamma$ et à $\partial E(\Gamma)$. Par le lemme $\underline{\overline{\mathrm{IX}}}$ on sait que l'angle $\alpha_1$ du secteur angulaire interieur (de $D_s$) determiné par $\partial E(\Gamma)$ est inferieur à $\pi$, \deleted{alors nécessairement $\alpha_1<\pi$}, car $s$, $s'$ et $s''$ ne sont pas colineaires. Par conséquent l'angle $\alpha_2$ qui sous-tend le secteur angulaire exterieur vérifie $\alpha_2>\pi$.

Mais $\operatorname{ext}\partial E(\Gamma)\subset\operatorname{ext}\Gamma$ implique qu'on peut trouver dans $D_s$ un angle $\alpha'_2>\pi$ qui sous-tend un secteur angulaire contenu dans $\operatorname{ext}\Gamma$, par conséquent le secteur angulaire interieur de $D_s$ \deleted{par} donné par $\Gamma$ a un angle $<\pi$.

Mais tout polygone convexe a au moins trois sommets et cela suffit à prouver le lemme.

On est en mésure de montrer le :

\uline{Lemme X}I : Si $\Gamma$ est un $J$-polygone \inserted{à $n\unsure{>|\geq}3$ \unsure{cotés|côtés}} il existe un couple de sommets $s,s'\in S(\Gamma)$ tels que $s'$ est visible de $s$ de l'interieur de $\Gamma$.
% Slanted pencil note to the left of the lemma heading.
\margin{faux pour $n=3$!}
% Pencil annotation below the lemma statement; its last two marks are dashes.
\margin{OPS $\Gamma$ non convexe (sinon c'est clair grâce à --)}

\uline{Démonstration} : Soit $s$ un sommet satisfaisant la propriété ennoncée au lemme $\underline{\overline{\mathrm{VIII}}}$%
% Dark-ink annotation at the left of the Roman reference.
\margin{\uline{lemme X}}, soient $s'$, $s''$ les deux sommets \deleted{consecutifs} \inserted{adjacents} à $s$ et $T$ le tri\uline{an}gle%
% Dark-ink note beside the underlined letters in triangle.
\margin{ambigu} de sommets $s,s',s''$. Si $(T-([s,s']\cup[s,s'']))\cap\Gamma=\unsure{\phi|\varnothing}$ on a que $]s',s''[\cap\Gamma=\phi$, donc $]s',s''[\subset\operatorname{int}\Gamma$ ou bien $]s',s''[\subset\operatorname{ext}\Gamma$. Mais $\inserted{]s',s''[\subset}T-([s,s']\cup[s,s''])\subset\deleted{\unsure{\operatorname{ext}}}\operatorname{int}\Gamma$ (puisqu'il est connexe et et son intersection avec $\operatorname{int}\Gamma$ est non vide par le choix de $s$) donc $]s',s''[\subset\operatorname{int}\Gamma$.
% The pencil leader curves rightward from the underlined ink note to the following pencil note.
\margin{\tikz[baseline=(fx.base)]{\node[inner sep=0pt] (fx) {\uline{faux}};\draw ([yshift=-.06cm]fx.south) .. controls ++(0,-.35) and ++(-.40,-.14) .. ++(1.3,-.12);}}
\margin{il faut prendre $s$ comme dans \unsure{le|la} \uline{dém} du lemme $\underline{\overline{\mathrm{VIII}}}$ (pas son énoncé!)}

% Slanted, heavily cancelled ink annotation immediately above the circled note.
\margin{\deleted{\illegible{2}}!}
% The detached diagonal mark before subset has a hooked lower end; it is not drawn across subset.
\uline{Si} \uline{$(T-([s,s']\cup[s,s'']))\cap\Gamma\mathrel{\text{\tikz[baseline=(rel.base)]{\node[inner sep=0pt] (rel) {$\subset$};\draw ([xshift=-.08cm,yshift=.30cm]rel.west)--([xshift=-.42cm,yshift=-.45cm]rel.west);\draw ([xshift=-.35cm,yshift=-.29cm]rel.west) to[bend right=15] ([xshift=-.22cm,yshift=-.41cm]rel.west);}}}]s',s''[$} on doit pouvoir trouver au moins un sommet $z\in]s',s''[$ (puisque $\Gamma\cap]s',s''[\ne]s',s''[$). On a aussi $\mathring{T}\cap\Gamma=\phi$ donc $]s,z[\subset\text{\tikz[baseline=(ic.base)]{\node[inner xsep=3pt,inner ysep=2pt] (ic) {$\mathring{T}\subset\operatorname{int}\Gamma$};\draw (ic.center) ellipse[x radius=1.02cm,y radius=.40cm];\draw ([yshift=-.40cm]ic.center) to[bend right=23] ++(.25,-.24);}}$ (puisque $\mathring{T}\subset\mathbb{R}^2-\Gamma$ et $\mathring{T}\cap\operatorname{int}\Gamma\ne\phi$)
% The initial mark has a long shaft and upturned right end; the separate forked oblique stroke beside the oval is also drawn.
\margin{\tikz[baseline=(m.base)]{\node[inner sep=1pt] (m) {\shortstack{$\unsure{\longrightarrow|\text{est}}$ non\\vide?}};\draw (m.center) ellipse[x radius=1.40cm,y radius=.50cm];\draw (-1.02,.72)--(-1.52,-.08);\draw (-1.40,.10)--(-1.34,-.25);}}
% Slanted pencil note linked to the oval around mathring T subset int Gamma.
\margin{dire avant! (et \unsure{tu as au lemme|tu as un lemme|tu auras lemme} $\underline{\overline{\mathrm{VIII}}}$)}
% archive sheet number: 35
